% basic nastaveni
\documentclass[12pt]{article}
\setlength{\parindent}{0pt}
\setlength{\parskip}{1em}
\usepackage[utf8]{inputenc}
\usepackage[czech]{babel}
\usepackage[T1]{fontenc}
\usepackage{graphics}
\usepackage{caption}
\usepackage{hyperref}
\usepackage{dsfont}
\usepackage{color}
\usepackage{float}
\usepackage{enumitem}
\usepackage{graphicx}

% matematicke package
\usepackage{tikz-cd}
\usetikzlibrary{babel}
\usepackage{amsmath}
\usepackage{amssymb}
\usepackage{amsfonts}
\usepackage{amssymb}
\usepackage{geometry}
\geometry{margin=2.5cm}
\usepackage{pgfplots}
\pgfplotsset{compat=1.18}
\usepackage{mathrsfs}

\title{Maturitní otázka číslo 20: Limita a spojitost funkce}
\author{}
\date{}

\begin{document}
\maketitle
Tato otázka se zabývá limitami, které zkoumají k jaké hodnotě se funkce blíží, když se její argument (x) nekonečně blíží k určité hodnotě. Je to nástroj pro analýzu bodů, kde funkce nemusí být definována. Spojitost zkoumá nepřerušovanost funkce, charakterizuje, kdy se funkce mění plynule bez "skoků". Funkce je spojitá, pokud její hodnota odpovídá tomu, co jsme předpověděli pomocí limity
\section*{Základní pojmy}
\subsection*{Okolí bodu}
Předpokládejme $a \in \mathbb{R}, \quad \epsilon \in \mathbb{R}>0$ 
Množinu $\{x\in\mathbb{R}; a-\epsilon<x<a+\epsilon\}$ nazýváme \textbf{okolí bodu} a. Číslo a nazýváme \textbf{střed okolí}, číslo $\epsilon$ nazýváme \textbf{poloměr okolí}.\\
Okolí bodu se značí jako $U(a,\epsilon)$
\begin{itemize}[itemsep = 2pt]
    \item \textbf{Pravé okolí bodu} nazýváme množinu $\{x\in \mathbb{R};a\leq x < a+\epsilon\}$ a značíme ho jako $U^+(a,\epsilon)$
    \item \textbf{Levé okolí bodu} nazýváme množinu $\{x \in \mathbb{R};a-\epsilon < x \leq a\}$ a značíme ho jako $U^-(a,\epsilon)$
    \item \textbf{Prstencové okolí bodu} nazýváme množinu $\{x\in \mathbb{R}; a-\epsilon < x < a+\epsilon \land x \neq a\}$ a značíme ho jako $P(a,\epsilon)$
    \item \textbf{Pravé prstencové okolí bodu} nazýváme množinu $\{x \in \mathbb{R};a<x<a+\epsilon\}$ a značíme ho jako $P^+(a,\epsilon)$
    \item \textbf{Levé prstencové okolí bodu} nazýváme množiny $\{x\in \mathbb{R}; a-\epsilon < x<a\}$ a značíme ho jako $P^-(a,\epsilon)$
    \item \textbf{Nevlastní body} nazýváme body $+\infty$ a $-\infty$.
\end{itemize}

Doplnit příklad

\section*{Spojitost funkce}
Řekněme, že funkce $f$ je spojitá v bodě $a$, jestliže \\
$\forall \epsilon > 0 \quad \exists \delta > 0 \quad \forall x \in \mathbb{R}\quad$ platí $\forall x, x \in U(a, \delta) \Rightarrow f(x) \in U(f(a), \epsilon)$ To znamená, že pokud pro všechna $\epsilon > 0$ existuje $\delta > 0$ a všechna $x \in \mathbb{R}$ tak platí, že všechna $x$ náleží do takového okolí $\delta$ bodu $a$, že funkční hodnota $x$ náleží do okolí $\epsilon$ funkční hodnoty $f(a)$

\textbf{Spojitost funkce zprava} nastává, pokud \\
$\forall\epsilon > 0 \quad \exists \delta>0 \quad \forall x \in \mathbb{R} $ platí $ \forall x \quad x\in U^+(a,\delta) \Rightarrow f(x) \in U(f(a), \epsilon)$ To znamená, že pokud pro všechna $\epsilon > 0$ existuje $\delta > 0$ a všechna $x \in \mathbb{R}$ tak platí, že všechna $x$ náleží do takového pravého okolí $\delta$ bodu $a$, že funkční hodnota $x$ náleží do okolí $\epsilon$ funkční hodnoty $f(a)$

\textbf{Spojitost funkce zleva} nastává, pokud\\ $\forall \epsilon > 0 \quad \exists \delta > 0 \quad \forall x \in \mathbb{R}$ platí $ \forall x \quad  x \in U^-(a,\delta) \Rightarrow f(x) \in U(f(a), \epsilon)$

\textbf{Definice spojitosti na otevřeném intervalu} $(a,b)$: Funkce je spojitá na tomto intervalu, pokud je spojitá v každém bodě tohoto intervalu
\textbf{Definice spojitosti na uzavřeném intervalu} $<a,b>$: Funkce je spojitá na tomto intervalu, je-li spojitá v $(a,b)$ a v bodě $a$ je spojitá zprava a v bodě $b$ je spojitá zleva

Příklad: Dokažte, že funkce $f(x) = x^2 - 2x$ je spojitá

Řešení: 

\subsection*{Věty o spojitosti}
Zeptat se Petrový(je jich strašně moc)

\subsection*{Přírůstek funkce v bodě}
Definice: Nechť je funkce $f(x)$ definována v jistém okolí bodu $x_0$. Nechť $\Delta x$ je libovolné reálné číslo, takové, že i bod $x_0 + \Delta x$ leží v definičním oboru funkce $f$.
Přírůstek funkce $f$ v bodě $x_0$ odpovídající přírůstku $\Delta x$ je hodnota:
\[\Delta f(x_0) = f(x_0 + \Delta x) - f(x_0)\]
Místo $\Delta x$ se používá i označení proměnnou $h$\\
Příklad: Zjistěte přírůstek funkce $f(x) = x^2$ v bodě $x = 2$
\[f(2 + h) = (2 + h)^2 = 4 + 4h + h^2\]
\[\Delta f = f(2 + h) - f(2)\]
\[\Delta f = 4 + 4h + h^2 - 4\]
\[\Delta f = 4h + h^2\]
Přírůstek v bodě $2$ této funkce je tedy $4h + h^2$, kde h je nějaký druhý bod, ke kterému se přírůstek měří. 



\end{document}